
%======================================================================
% 附录A：矢量与张量公式汇总
%======================================================================
\chapter*{附录A：矢量与张量公式汇总}
\addcontentsline{toc}{chapter}{附录A：矢量与张量公式汇总}

本章汇总等离子体物理中频繁使用的矢量分析与张量运算公式，方便查阅。所有公式在 SI 单位制下给出。

\section{直角坐标系 $(x,y,z)$}

设标量场为 $\phi(x,y,z)$，矢量场为 $\vect{A}=A_x\unitvec{x}+A_y\unitvec{y}+A_z\unitvec{z}$。

\begin{definition}{梯度}{}
\begin{equation}
\grad\phi = \frac{\partial \phi}{\partial x}\unitvec{x} + \frac{\partial \phi}{\partial y}\unitvec{y} + \frac{\partial \phi}{\partial z}\unitvec{z}
\end{equation}
\end{definition}

\begin{definition}{散度}{}
\begin{equation}
\divg\vect{A} = \frac{\partial A_x}{\partial x} + \frac{\partial A_y}{\partial y} + \frac{\partial A_z}{\partial z}
\end{equation}
\end{definition}

\begin{definition}{旋度}{}
\begin{equation}
\curl\vect{A} =
\begin{vmatrix}
\unitvec{x} & \unitvec{y} & \unitvec{z} \\[4pt]
\frac{\partial }{\partial x} & \frac{\partial }{\partial y} & \frac{\partial }{\partial z} \\[4pt]
A_x & A_y & A_z
\end{vmatrix}
= \left(\frac{\partial A_z}{\partial y}-\frac{\partial A_y}{\partial z}\right)\unitvec{x}
+ \left(\frac{\partial A_x}{\partial z}-\frac{\partial A_z}{\partial x}\right)\unitvec{y}
+ \left(\frac{\partial A_y}{\partial x}-\frac{\partial A_x}{\partial y}\right)\unitvec{z}
\end{equation}
\end{definition}

\begin{definition}{拉普拉斯算子}{}
对标量场：
\begin{equation}
\lap\phi = \frac{\partial^2 \phi}{\partial x^2} + \frac{\partial^2 \phi}{\partial x^2} + \frac{\partial^2 \phi}{\partial x^2}
\end{equation}
对矢量场：
\begin{equation}
\lap\vect{A} = \lap A_x\,\unitvec{x} + \lap A_y\,\unitvec{y} + \lap A_z\,\unitvec{z}
\end{equation}
\end{definition}


\section{柱坐标系 $(r,\theta,z)$}

设标量场为 $\phi(r,\theta,z)$，矢量场为 $\vect{A}=A_r\unitvec{r}+A_\theta\unitvec{\theta}+A_z\unitvec{z}$。

\begin{definition}{梯度}{}
\begin{equation}
\grad\phi = \frac{\partial \phi}{\partial r}\unitvec{r} + \frac{1}{r}\frac{\partial \phi}{\partial \theta}\unitvec{\theta} + \frac{\partial \phi}{\partial z}\unitvec{z}
\end{equation}
\end{definition}

\begin{definition}{散度}{}
\begin{equation}
\divg\vect{A} = \frac{1}{r}\frac{\partial }{\partial r}(rA_r) + \frac{1}{r}\frac{\partial A_\theta}{\partial \theta} + \frac{\partial A_z}{\partial z}
\end{equation}
\end{definition}

\begin{definition}{旋度}{}
\begin{equation}
\curl\vect{A} =
\left(\frac{1}{r}\frac{\partial A_z}{\partial \theta}-\frac{\partial A_\theta}{\partial z}\right)\unitvec{r}
+ \left(\frac{\partial A_r}{\partial z}-\frac{\partial A_z}{\partial r}\right)\unitvec{\theta}
+ \left(\frac{1}{r}\frac{\partial }{\partial r}(rA_\theta)-\frac{1}{r}\frac{\partial A_r}{\partial \theta}\right)\unitvec{z}
\end{equation}
\end{definition}

\begin{definition}{拉普拉斯算子}{}
\begin{equation}
\lap\phi = \frac{1}{r}\frac{\partial }{\partial r}\!\left(r\frac{\partial \phi}{\partial r}\right)
+ \frac{1}{r^2}\frac{\partial^2 \phi}{\partial x^2} + \frac{\partial^2 \phi}{\partial x^2}
\end{equation}
\end{definition}

\begin{insight}[轴对称假设下的简化]
在等离子体柱（如托卡马克）中常取轴对称 $\partial/\partial\theta=0$。此时柱坐标旋度的 $\theta$ 分量和散度均显著简化，许多推导可在一维或二维径向坐标中完成。
\end{insight}


\section{球坐标系 $(r,\theta,\varphi)$}

设标量场为 $\phi(r,\theta,\varphi)$，矢量场为 $\vect{A}=A_r\unitvec{r}+A_\theta\unitvec{\theta}+A_\varphi\unitvec{\varphi}$。以下公式在恒星大气、日球层物理中频繁使用。

\begin{definition}{梯度}{}
\begin{equation}
\grad\phi = \frac{\partial \phi}{\partial r}\unitvec{r}
+ \frac{1}{r}\frac{\partial \phi}{\partial \theta}\unitvec{\theta}
+ \frac{1}{r\sin\theta}\frac{\partial \phi}{\partial \varphi}\unitvec{\varphi}
\end{equation}
\end{definition}

\begin{definition}{散度}{}
\begin{equation}
\divg\vect{A} = \frac{1}{r^2}\frac{\partial }{\partial r}(r^2A_r)
+ \frac{1}{r\sin\theta}\frac{\partial }{\partial \theta}(\sin\theta\,A_\theta)
+ \frac{1}{r\sin\theta}\frac{\partial A_\varphi}{\partial \varphi}
\end{equation}
\end{definition}

\begin{definition}{旋度}{}
\begin{multline}
\curl\vect{A} =
\frac{1}{r\sin\theta}\left[\frac{\partial }{\partial \theta}(\sin\theta\,A_\varphi)-\frac{\partial A_\theta}{\partial \varphi}\right]\unitvec{r}
+ \frac{1}{r}\left[\frac{1}{\sin\theta}\frac{\partial A_r}{\partial \varphi}-\frac{\partial }{\partial r}(rA_\varphi)\right]\unitvec{\theta} \\[6pt]
+ \frac{1}{r}\left[\frac{\partial }{\partial r}(rA_\theta)-\frac{\partial A_r}{\partial \theta}\right]\unitvec{\varphi}
\end{multline}
\end{definition}

\begin{definition}{拉普拉斯算子}{}
\begin{equation}
\lap\phi = \frac{1}{r^2}\frac{\partial }{\partial r}\!\left(r^2\frac{\partial \phi}{\partial r}\right)
+ \frac{1}{r^2\sin\theta}\frac{\partial }{\partial \theta}\!\left(\sin\theta\frac{\partial \phi}{\partial \theta}\right)
+ \frac{1}{r^2\sin^2\theta}\frac{\partial^2 \phi}{\partial x^2}
\end{equation}
\end{definition}


\section{常用矢量恒等式}

以下恒等式对任意足够光滑的矢量场 $\vect{A},\vect{B},\vect{C}$ 和标量场 $\phi$ 成立。

\begin{theorem}{基本恒等式}{}
\begin{align}
\vect{A}\times(\vect{B}\times\vect{C}) &= (\vect{A}\cdot\vect{C})\vect{B} - (\vect{A}\cdot\vect{B})\vect{C} \label{eq:bac-cab} \\[6pt]
\vect{A}\cdot(\vect{B}\times\vect{C}) &= \vect{B}\cdot(\vect{C}\times\vect{A}) = \vect{C}\cdot(\vect{A}\times\vect{B}) \\[6pt]
(\vect{A}\times\vect{B})\cdot(\vect{C}\times\vect{D}) &= (\vect{A}\cdot\vect{C})(\vect{B}\cdot\vect{D}) - (\vect{A}\cdot\vect{D})(\vect{B}\cdot\vect{C})
\end{align}
\end{theorem}

\begin{theorem}{含 $\nabla$ 的微分恒等式}{}
\begin{align}
\grad(\vect{A}\cdot\vect{B}) &= (\vect{A}\cdot\grad)\vect{B} + (\vect{B}\cdot\grad)\vect{A} + \vect{A}\times(\curl\vect{B}) + \vect{B}\times(\curl\vect{A}) \\[6pt]
\divg(\vect{A}\times\vect{B}) &= \vect{B}\cdot(\curl\vect{A}) - \vect{A}\cdot(\curl\vect{B}) \\[6pt]
\curl(\vect{A}\times\vect{B}) &= \vect{A}(\divg\vect{B}) - \vect{B}(\divg\vect{A}) + (\vect{B}\cdot\grad)\vect{A} - (\vect{A}\cdot\grad)\vect{B} \\[6pt]
\curl(\grad\phi) &= \vect{0} \\[6pt]
\divg(\curl\vect{A}) &= 0 \\[6pt]
\curl(\curl\vect{A}) &= \grad(\divg\vect{A}) - \lap\vect{A} \label{eq:curlcurl}
\end{align}
\end{theorem}

式 \eqref{eq:curlcurl} 在推导 MHD 中的涡量方程和阿尔芬波色散关系时至关重要。


\section{极简二阶张量入门}

等离子体物理中，MHD 动量方程涉及应力张量，这里给出仅够本书使用的最低限度张量知识。

\begin{definition}{并矢}{}
两个矢量 $\vect{A}$ 与 $\vect{B}$ 的并矢（dyad）记为 $\vect{A}\vect{B}$（或 $\vect{A}\otimes\vect{B}$），其分量为 $(\vect{A}\vect{B})_{ij}=A_iB_j$。在直角坐标下可写为矩阵：
\begin{equation}
\vect{A}\vect{B} =
\begin{pmatrix}
A_xB_x & A_xB_y & A_xB_z \\
A_yB_x & A_yB_y & A_yB_z \\
A_zB_x & A_zB_y & A_zB_z
\end{pmatrix}
\end{equation}
\end{definition}

\begin{definition}{张量散度}{}
对二阶张量 $\tens{T}$，其散度 $\divg\tens{T}$ 是一个矢量，分量定义为
\begin{equation}
(\divg\tens{T})_i = \sum_j \frac{\partial T_{ij}}{\partial x_j}
\end{equation}
\end{definition}

\begin{theorem}{MHD 磁应力张量}{}
磁应力张量（Maxwell 应力张量的磁场部分）为
\begin{equation}
\tens{T}_M = \frac{1}{\mu_0}\left(\vect{B}\vect{B} - \frac{B^2}{2}\tens{I}\right)
\end{equation}
其中 $\tens{I}$ 为单位张量。其散度给出磁场对等离子体的作用力：
\begin{equation}
\divg\tens{T}_M = \frac{1}{\mu_0}(\curl\vect{B})\times\vect{B}
= \frac{1}{\mu_0}\left[ (\vect{B}\cdot\grad)\vect{B} - \grad\left(\frac{B^2}{2}\right) \right]
\end{equation}
\end{theorem}

\begin{tip}
记忆口诀：磁应力张量可看作各向异性压力——沿磁场方向为张力 $B^2/\mu_0$，垂直方向为压力 $B^2/(2\mu_0)$。在 MHD 平衡分析中，这一直觉比分量公式更常用。
\end{tip}


%======================================================================
% 附录B：等离子体特征参数速查表
%======================================================================
\chapter*{附录B：等离子体特征参数速查表}
\addcontentsline{toc}{chapter}{附录B：等离子体特征参数速查表}

本章将书中所有关键等离子体参数汇总为一张速查表，按物理内涵分组，便于检索与对照。


\section{参数总表}

\begin{table}[htbp]
\centering
\caption{等离子体特征参数速查表（SI 单位制，$T$ 以 eV 为单位时，$k_\mathrm{B}T\to T$）}
\label{tab:plasma-parameters}
\small
\begin{tabular}{@{}p{2.8cm}p{1.5cm}p{6.5cm}p{3.5cm}p{1.8cm}@{}}
\toprule
\textbf{参数名} & \textbf{符号} & \textbf{公式} & \textbf{典型数值} & \textbf{章节} \\
\midrule
\multicolumn{5}{l}{\textit{一、基本定义}} \\
\midrule
德拜长度 & $\lambda_D$ & $\displaystyle\lambda_D = \sqrt{\frac{\eps_0 k_\mathrm{B} T_e}{n_e e^2}}$ & $10^{-4}\sim10^{-3}$ m (聚变) & \S 2.2 \\
等离子体频率 & $\omega_{pe}$ & $\displaystyle\omega_{pe} = \sqrt{\frac{n_e e^2}{\eps_0 m_e}}$ & $10^{9}\sim10^{12}$ s$^{-1}$ & \S 2.3 \\
离子等离子体频率 & $\omega_{pi}$ & $\displaystyle\omega_{pi} = \sqrt{\frac{n_i Z^2 e^2}{\eps_0 m_i}}$ & $\omega_{pi}\ll\omega_{pe}$ & \S 2.3 \\
\midrule
\multicolumn{5}{l}{\textit{二、静电相关}} \\
\midrule
热速度（电子） & $v_{te}$ & $\displaystyle v_{te} = \sqrt{\frac{k_\mathrm{B}T_e}{m_e}}$ & $10^{6}\sim10^{7}$ m/s & \S 3.1 \\
热速度（离子） & $v_{ti}$ & $\displaystyle v_{ti} = \sqrt{\frac{k_\mathrm{B}T_i}{m_i}}$ & $10^{4}\sim10^{5}$ m/s & \S 3.1 \\
德拜球粒子数 & $N_D$ & $\displaystyle N_D = n\lambda_D^3$ & $\gg 1$（理想等离子体）& \S 2.2 \\
耦合参数 & $\Gamma$ & $\displaystyle\Gamma = \frac{e^2}{4\pi\eps_0 a k_\mathrm{B}T}$，$a=(3/4\pi n)^{1/3}$ & $\ll 1$（弱耦合）& \S 2.5 \\
\midrule
\multicolumn{5}{l}{\textit{三、磁化相关}} \\
\midrule
电子回旋频率 & $\omega_{ce}$ & $\displaystyle\omega_{ce} = \frac{|e|B}{m_e}$ & $10^{10}\sim10^{12}$ s$^{-1}$ & \S 3.2 \\
离子回旋频率 & $\omega_{ci}$ & $\displaystyle\omega_{ci} = \frac{|Ze|B}{m_i}$ & $10^{6}\sim10^{8}$ s$^{-1}$ & \S 3.2 \\
电子回旋半径 & $r_{ce}$ & $\displaystyle r_{ce} = \frac{v_{\perp e}}{\omega_{ce}} = \frac{m_e v_{\perp e}}{|e|B}$ & $10^{-5}\sim10^{-3}$ m & \S 3.2 \\
离子回旋半径 & $r_{ci}$ & $\displaystyle r_{ci} = \frac{v_{\perp i}}{\omega_{ci}} = \frac{m_i v_{\perp i}}{|Ze|B}$ & $10^{-3}\sim10^{-1}$ m & \S 3.2 \\
阿尔芬速度 & $v_A$ & $\displaystyle v_A = \frac{B}{\sqrt{\mu_0\rho_m}}$，$\rho_m=n_i m_i$ & $10^{5}\sim10^{7}$ m/s & \S 7.3 \\
离子声速 & $c_s$ & $\displaystyle c_s = \sqrt{\frac{\gamma_e Z k_\mathrm{B}T_e + \gamma_i k_\mathrm{B}T_i}{m_i}}$ & $10^{4}\sim10^{5}$ m/s & \S 8.2 \\
安全因子 & $q$ & $\displaystyle q = \frac{r B_\varphi}{R B_\theta}$（托卡马克）& $q>1$（稳定）& \S 11.2 \\
比压 & $\beta$ & $\displaystyle\beta = \frac{2\mu_0 p}{B^2} = \frac{2\mu_0 n k_\mathrm{B}(T_e+T_i)}{B^2}$ & $0.01\sim0.1$（托卡马克）& \S 11.3 \\
\midrule
\multicolumn{5}{l}{\textit{四、碰撞输运}} \\
\midrule
朗道长度 & $b_\pi$ & $\displaystyle b_\pi = \frac{e^2}{4\pi\eps_0 k_\mathrm{B}T}$ & $\sim 10^{-14}$ m & \S 4.1 \\
库仑对数 & $\ln\Lambda$ & $\displaystyle\ln\Lambda = \ln\!\left(\frac{\lambda_D}{b_\pi}\right)$ & $10\sim 25$ & \S 4.1 \\
电子--离子碰撞频率 & $\nu_{ei}$ & $\displaystyle\nu_{ei} = \frac{n_i Z^2 e^4}{3(2\pi)^{3/2}\eps_0^2 m_e^{1/2}(k_\mathrm{B}T_e)^{3/2}}\ln\Lambda$ & $\sim 10^{5}$ s$^{-1}$ & \S 4.2 \\
Spitzer 电阻率 & $\eta_\parallel$ & $\displaystyle\eta_\parallel = \frac{Z e^2 m_e^{1/2}}{3(2\pi)^{3/2}\eps_0^2(k_\mathrm{B}T_e)^{3/2}}\ln\Lambda$ & $\sim 10^{-8}$ $\Omega\cdot$m & \S 4.3 \\
双极扩散系数 & $D_a$ & $\displaystyle D_a \approx \frac{\mu_e D_i + \mu_i D_e}{\mu_e + \mu_i}\approx D_i\left(1+\frac{T_e}{T_i}\right)$ & $\sim D_i$（被 $T_e/T_i$ 放大）& \S 6.4 \\
\midrule
\multicolumn{5}{l}{\textit{五、无量纲参数}} \\
\midrule
磁雷诺数 & $R_m$ & $\displaystyle R_m = \frac{\mu_0 v L}{\eta}$ & $\gg 1$（理想 MHD）& \S 9.1 \\
等离子体 beta & $\beta$ & 同上文定义 & 见上文 & \S 11.3 \\
\bottomrule
\end{tabular}
\end{table}

\begin{insight}{}{}
表 \ref{tab:plasma-parameters} 中的公式均使用 SI 单位。若采用高斯单位，需将 $1/(4\pi\eps_0)\to 1$，$\mu_0\to 4\pi/c^2$，并注意 $B$（高斯）$=10^4 B$（特斯拉）。
\end{insight}

\begin{warning}[温度单位的陷阱]
本书中公式若写 $T_e$（未标注 $k_\mathrm{B}$），默认以 \textbf{eV} 为温度单位。此时能量公式 $k_\mathrm{B}T_e\to T_e$ 仍成立，因为 1 eV 在数值上等于 $k_\mathrm{B}\times 11604$ K。但在做量纲分析时，必须显式写回 $k_\mathrm{B}$。
\end{warning}


%======================================================================
% 附录C：书中省略的中间数学步骤
%======================================================================
\chapter*{附录C：书中省略的中间数学步骤}
\addcontentsline{toc}{chapter}{附录C：书中省略的中间数学步骤}

正文中为保持叙述流畅，省略了若干冗长但重要的数学推导。本章将它们完整补出，并标注对应正文章节。如果读者在正文处感到跳跃，可查阅本节。


\section{C.1 广义欧姆定律的完整推导}
\label{sec:gen-ohm}

\textbf{对应正文：}\S 5.2（双流体方程到 MHD 的过渡）。

从电子和离子的动量方程出发：
\begin{align}
m_i n_i \left(\frac{\partial \vect{v}_i}{\partial t} + (\vect{v}_i\cdot\grad)\vect{v}_i\right) &= n_i e(\vect{E}+\vect{v}_i\times\vect{B}) - \grad p_i - \vect{R} \label{eq:ion-mom}\\
m_e n_e \left(\frac{\partial \vect{v}_e}{\partial t} + (\vect{v}_e\cdot\grad)\vect{v}_e\right) &= -n_e e(\vect{E}+\vect{v}_e\times\vect{B}) - \grad p_e + \vect{R} \label{eq:elec-mom}
\end{align}
其中 $\vect{R}=m_e n_e \nu_{ei}(\vect{v}_i-\vect{v}_e)$ 为摩擦项（碰撞动量交换）。定义总质量密度 $\rho_m=m_i n_i+m_e n_e\approx m_i n_i$、质心速度 $\vect{v}=(m_i n_i\vect{v}_i+m_e n_e\vect{v}_e)/\rho_m$，以及电流密度 $\vect{j}=n_i e\vect{v}_i-n_e e\vect{v}_e$。

\textbf{步骤 1：} 将 \eqref{eq:ion-mom} 与 \eqref{eq:elec-mom} 相加，消去 $\vect{R}$，得到总动量方程（即 MHD 动量方程）：
\begin{equation}
\rho_m \left(\frac{\partial \vect{v}}{\partial t} + (\vect{v}\cdot\grad)\vect{v}\right) = \vect{j}\times\vect{B} - \grad p
\end{equation}
其中 $p=p_i+p_e$。

\textbf{步骤 2：} 将 \eqref{eq:elec-mom} 除以 $-n_e e$，并用 $\vect{v}_e = \vect{v}_i - \vect{j}/(n_e e)$ 消去 $\vect{v}_e$，整理得：
\begin{equation}
\vect{E} + \vect{v}_i\times\vect{B} = \frac{1}{n_e e}\vect{R} + \frac{m_e}{n_e e^2}\frac{\partial \vect{j}}{\partial t} + \frac{1}{n_e e}\grad p_e - \frac{1}{n_e e}\vect{j}\times\vect{B}
\end{equation}

\textbf{步骤 3：} 代入 $\vect{R}=m_e n_e \nu_{ei}(\vect{v}_i-\vect{v}_e)=m_e \nu_{ei}\vect{j}/e$，并忽略电子惯性项（$m_e\to 0$），令 $\vect{v}_i\approx\vect{v}$，得到广义欧姆定律：
\begin{equation}
\boxed{
\vect{E} + \vect{v}\times\vect{B} = \eta\vect{j} + \frac{1}{n_e e}\grad p_e - \frac{1}{n_e e}\vect{j}\times\vect{B}
}
\end{equation}
其中 $\eta=m_e\nu_{ei}/(n_e e^2)$ 为电阻率。右边三项分别对应：电阻、电子压力梯度（双极扩散效应）、霍尔效应。


\section{C.2 Landau 围道积分的详细推导}
\label{sec:landau-contour}

\textbf{对应正文：}\S 8.4（无碰撞阻尼）。\textbf{难度标注：}[挑战]

Vlasov--Poisson 体系线性化后，电子密度扰动的响应由积分
\begin{equation}
\int_{-\infty}^{+\infty} \frac{\partial F_0/\partial v}{v-\omega/k}\,\diff v
\end{equation}
给出。当 $k>0$ 时，若波的相位速度 $\omega/k$ 位于实轴上，积分在 $v=\omega/k$ 处出现奇点。Landau 的处理是：令 $\omega\to\omega+\ii\gamma$（$\gamma>0$ 为微小增长率），则奇点移入复平面上半平面；按因果律，积分路径应沿实轴从下方绕过奇点。取 $\gamma\to 0^+$ 的极限，得到主值积分加半留数：
\begin{equation}
\int_{-\infty}^{+\infty} \frac{\partial F_0/\partial v}{v-\omega/k}\,\diff v
= \mathcal{P}\int_{-\infty}^{+\infty} \frac{\partial F_0/\partial v}{v-\omega/k}\,\diff v
+ \ii\pi\left.\frac{\partial F_0}{\partial v}\right|_{v=\omega/k}
\end{equation}

对 Maxwell 分布 $F_0(v)\propto\exp(-v^2/v_{te}^2)$，有
\begin{equation}
\left.\frac{\partial F_0}{\partial v}\right|_{v=\omega/k} < 0 \quad (\omega/k>0)
\end{equation}
因此虚部贡献为负，对应阻尼率 $\gamma<0$。

\begin{insight}{}{}
Landau 围道的物理本质是 \textbf{因果律}：在 $t=-\infty$ 时扰动为零，这要求积分围道必须使上半平面的极点贡献被包含。如果反过来取 $\gamma<0$，将得到违反因果律的增长解。
\end{insight}


\section{C.3 MHD 阿尔芬波色散关系的完整推导}
\label{sec:alfven-dispersion}

\textbf{对应正文：}\S 7.3（MHD 波）。

从理想 MHD 线性化方程组（均匀背景 $B_0\unitvec{z}$，$\rho_0, p_0$ 为常数）：
\begin{align}
\ii\omega\rho_0\vect{v}_1 &= \frac{1}{\mu_0}(\curl\vect{B}_1)\times\vect{B}_0 - \grad p_1 \label{eq:lin-mom}\\
\vect{B}_1 &= -\frac{1}{\ii\omega}\curl(\vect{v}_1\times\vect{B}_0) \label{eq:lin-far}\\
\ii\omega p_1 &= -\gamma p_0\divg\vect{v}_1 \label{eq:lin-eos}
\end{align}

取傅里叶模式 $\exp[\ii(k_z z - \omega t)]$，设 $\vect{v}_1$ 沿 $x$ 方向且仅依赖于 $z$（即剪切阿尔芬波，$\vect{k}\parallel\vect{B}_0$，$\vect{v}_1\perp\vect{B}_0$，$k_\perp=0$）。

\textbf{步骤 1：} 由 \eqref{eq:lin-far}，
\begin{equation}
\vect{B}_1 = -\frac{1}{\ii\omega}\curl(v_1\unitvec{x}\times B_0\unitvec{z}) = -\frac{B_0 v_1}{\ii\omega}\curl(-\unitvec{y}) = -\frac{B_0 k_z v_1}{\omega}\unitvec{x}
\end{equation}

\textbf{步骤 2：} 将上式代入 \eqref{eq:lin-mom}。注意 $(\curl\vect{B}_1)\times\vect{B}_0$ 中，$\vect{B}_1$ 沿 $x$，故 $\curl\vect{B}_1$ 沿 $y$，再叉乘 $B_0\unitvec{z}$ 回到 $x$：
\begin{equation}
\ii\omega\rho_0 v_1 = \frac{1}{\mu_0}(-\ii k_z B_{1x})(-B_0) - 0 = \frac{\ii k_z B_0^2 k_z v_1}{\mu_0\omega}
\end{equation}

\textbf{步骤 3：} 整理得 $\omega^2 = k_z^2 B_0^2/(\mu_0\rho_0) = k_z^2 v_A^2$，即
\begin{equation}
\boxed{\omega = \pm k_z v_A}
\end{equation}

此即剪切阿尔芬波的色散关系。对压缩阿尔芬波（快磁声波），需保留 $k_\perp\neq 0$ 和 $p_1$ 项，推导见正文 \S 7.3。


\section{C.4 Grad--Shafranov 方程的简化推导}
\label{sec:grad-shafranov}

\textbf{对应正文：}\S 11.2（MHD 平衡）。

在轴对称托卡马克中 $(R,\varphi,Z)$，磁场可写为
\begin{equation}
\vect{B} = \frac{1}{R}\grad\psi\times\unitvec{\varphi} + \frac{F(\psi)}{R}\unitvec{\varphi}
\end{equation}
其中 $\psi(R,Z)$ 为极向磁通函数，$F(\psi)=RB_\varphi$ 为环向磁场相关的流函数。

\textbf{步骤 1：} 由 $\vect{j}\times\vect{B}=\grad p$ 的 $\varphi$ 分量，得 $p=p(\psi)$。再由 $\divg\vect{B}=0$ 自动满足。

\textbf{步骤 2：} 将 $\vect{j}=(1/\mu_0)\curl\vect{B}$ 的极向分量与力的平衡投影到 $\grad\psi$ 方向，得到：
\begin{equation}
\nabla^{2*}\psi = -\mu_0 R^2 \frac{\diff p}{\diff\psi} - F\frac{\diff F}{\diff\psi}
\end{equation}
其中椭圆算子（Grad--Shafranov 算子）定义为
\begin{equation}
\nabla^{2*}\psi \equiv R^2\divg\left(\frac{1}{R^2}\grad\psi\right) = R\frac{\partial }{\partial R}\!\left(\frac{1}{R}\frac{\partial \psi}{\partial R}\right) + \frac{\partial^2 \phi}{\partial x^2}
\end{equation}

\textbf{步骤 3：} 给定 $p(\psi)$ 和 $F(\psi)$（两个自由剖面函数），在边界条件 $\psi=\psi_b$（导体壁）下求解椭圆方程，即得托卡马克平衡位形。

\begin{tip}
Grad--Shafranov 方程是二维椭圆型偏微分方程，其解析解仅对特定的 $p(\psi)$ 和 $F(\psi)$ 选择才存在（如 Solov'ev 解）。数值上通常使用有限元或迭代法（如 Picard 迭代）求解。
\end{tip}


%======================================================================
% 附录D：单位换算与物理常数
%======================================================================
\chapter*{附录D：单位换算与物理常数}
\addcontentsline{toc}{chapter}{附录D：单位换算与物理常数}

本章提供等离子体物理中常用的物理常数、单位换算关系，以及 SI 与高斯单位制之间的对照。


\section{基础物理常数}

\begin{table}[htbp]
\centering
\caption{基础物理常数（精确到 4 位有效数字）}
\label{tab:constants}
\begin{tabular}{@{}llll@{}}
\toprule
\textbf{常数名} & \textbf{符号} & \textbf{数值} & \textbf{单位} \\
\midrule
真空光速 & $c$ & $2.998\times 10^{8}$ & m/s \\
真空磁导率 & $\mu_0$ & $4\pi\times 10^{-7} = 1.257\times 10^{-6}$ & H/m \\
真空介电常数 & $\eps_0$ & $8.854\times 10^{-12}$ & F/m \\
元电荷 & $e$ & $1.602\times 10^{-19}$ & C \\
电子质量 & $m_e$ & $9.109\times 10^{-31}$ & kg \\
质子质量 & $m_p$ & $1.673\times 10^{-27}$ & kg \\
中子质量 & $m_n$ & $1.675\times 10^{-27}$ & kg \\
玻尔兹曼常数 & $k_\mathrm{B}$ & $1.381\times 10^{-23}$ & J/K \\
普朗克常数 & $h$ & $6.626\times 10^{-34}$ & J$\cdot$s \\
约化普朗克常数 & $\hbar$ & $1.055\times 10^{-34}$ & J$\cdot$s \\
精细结构常数 & $\alpha$ & $7.297\times 10^{-3}$ & -- \\
阿伏伽德罗常数 & $N_A$ & $6.022\times 10^{23}$ & mol$^{-1}$ \\
理想气体常数 & $R$ & $8.314$ & J/(mol$\cdot$K) \\
\bottomrule
\end{tabular}
\end{table}


\section{SI 与高斯单位制换算}

等离子体物理文献中，高斯（CGS）单位仍广泛使用。以下给出关键电磁量的公式对照。

\begin{table}[htbp]
\centering
\caption{SI 与高斯单位制对照表}
\label{tab:unit-conversion}
\begin{tabular}{@{}p{3.5cm}p{5.5cm}p{5.5cm}@{}}
\toprule
\textbf{物理量} & \textbf{SI 公式} & \textbf{高斯（CGS）公式} \\
\midrule
电场强度 & $\vect{E}$ [V/m] & $\vect{E}$ [statV/cm] $= \vect{E}_\text{SI}/(3\times 10^4)$ \\
磁感应强度 & $\vect{B}$ [T] $=\mu_0(\vect{H}+\vect{M})$ & $\vect{B}$ [G] $= 10^4\vect{B}_\text{SI}$ \\
电荷密度 & $\rho = nq$ [C/m$^3$] & $\rho$ [statC/cm$^3$] $= \rho_\text{SI}/(3\times 10^9)$ \\
电流密度 & $\vect{j}$ [A/m$^2$] & $\vect{j}$ [statA/cm$^2$] $= \vect{j}_\text{SI}/(3\times 10^5)$ \\
洛伦兹力 & $\vect{F}=q(\vect{E}+\vect{v}\times\vect{B})$ & $\vect{F}=q\left(\vect{E}+\frac{\vect{v}}{c}\times\vect{B}\right)$ \\
麦克斯韦（安培） & $\curl\vect{B}=\mu_0\vect{j}+\mu_0\eps_0\partial\vect{E}/\partial t$ & $\curl\vect{B}=\frac{4\pi}{c}\vect{j}+\frac{1}{c}\partial\vect{E}/\partial t$ \\
库仑力 & $F=\frac{1}{4\pi\eps_0}\frac{q_1q_2}{r^2}$ & $F=\frac{q_1q_2}{r^2}$ \\
\bottomrule
\end{tabular}
\end{table}

\begin{warning}[磁感应强度的称呼]
SI 中 $\vect{B}$ 称为磁感应强度（特斯拉），$\vect{H}$ 为磁场强度（A/m）；高斯单位中通常不区分，统称磁场 $\vect{B}$（高斯）。$1\,\text{T} = 10^4\,\text{G}$。本书 SI 公式中一律使用 $\vect{B}$。
\end{warning}


\section{常用能量、密度、压强换算}

\begin{table}[htbp]
\centering
\caption{常用换算关系}
\label{tab:common-conversion}
\begin{tabular}{@{}lll@{}}
\toprule
\textbf{换算} & \textbf{关系} & \textbf{备注} \\
\midrule
温度 $T$ & $1\,\text{eV} = 11604\,\text{K} = 1.602\times 10^{-19}\,\text{J}$ & 等离子体中常用 eV 作温度单位 \\
数密度 $n$ & $1\,\text{cm}^{-3} = 10^6\,\text{m}^{-3}$ & 空间物理常用 cm$^{-3}$ \\
磁场 $B$ & $1\,\text{T} = 10^4\,\text{G}$ & 聚变物理常用 T，天体物理常用 G \\
压强 $p$ & $1\,\text{Pa} = 10\,\text{dyn/cm}^2 = 7.501\,\text{mTorr}$ & 真空物理常用 Torr/mTorr \\
& $1\,\text{atm} = 1.013\times 10^5\,\text{Pa} = 760\,\text{Torr}$ & \\
\bottomrule
\end{tabular}
\end{table}

\begin{tip}
在等离子体物理中，一个极其有用的经验公式是：当 $T_e$ 以 eV 为单位时，
\begin{equation}
v_{te} [\text{m/s}] \approx 4.19\times 10^5\sqrt{T_e [\text{eV}]}
\end{equation}
此式来自 $v_{te}=\sqrt{k_\mathrm{B}T_e/m_e}$，将 eV 换算为焦耳后代入常数即得。
\end{tip}


%======================================================================
% 附录E：分级自学书单
%======================================================================
\chapter*{附录E：分级自学书单}
\addcontentsline{toc}{chapter}{附录E：分级自学书单}

本章为不同阶段的读者提供等离子体物理学习路径。从入门到前沿，循序渐进。


\section{入门级（中文优先）}

适合零基础或仅具备普通物理背景的读者，中文教材可降低语言门槛，专注物理概念。

\begin{enumerate}
\item \textbf{金晓峰，《等离子体物理导论》}，北京大学出版社。\\
面向物理系本科高年级，概念清晰，公式推导详尽，对国内学生非常友好。重点章节：第1--5章（基础概念、单粒子运动、双流体、MHD）。

\item \textbf{郑春开，《等离子体物理学》}，北京大学出版社。\\
体系完整，涵盖基础到磁约束聚变。特点是 MHD 平衡与稳定性章节写得尤为出色，适合作为系统学习的主干教材。
\end{enumerate}


\section{进阶级（英文教材）}

具备一定基础后，建议直接阅读英文教材，术语与前沿文献接轨。以下三本互补性极强。

\begin{enumerate}
\item \textbf{Francis F. Chen, \textit{Introduction to Plasma Physics and Controlled Fusion}, 3rd ed., Springer.}\\
最经典的入门教材，语言极其通俗。第1卷（等离子体物理）覆盖单粒子运动、动力学、MHD、波；第2卷（受控聚变）涉及托卡马克、仿星器。每一章都有大量例题和数值估计，非常适合自学者。

\item \textbf{Robert J. Goldston \& Paul H. Rutherford, \textit{Introduction to Plasma Physics}, IOP Publishing.}\\
普林斯顿等离子体物理实验室（PPPL）的授课教材。相比 Chen，此书对碰撞输运、波--粒子相互作用、MHD 稳定性有更深入的处理，数值例子偏聚变实验。

\item \textbf{Richard Fitzpatrick, \textit{Plasma Physics: An Introduction}, CRC Press.}\\
在线免费版（https://farside.ph.utexas.edu/plasma）可用。该书从动力学理论出发，对 Landau 阻尼、碰撞算符、MHD 波的推导极为严谨，适合希望补全数学细节的读者。
\end{enumerate}


\section{专业级（英文专著）}

进入科研阶段后，需要针对特定方向深入钻研。

\begin{enumerate}
\item \textbf{Per Helander \& Dieter J. Sigmar, \textit{Collisional Transport in Magnetized Plasmas}, Cambridge.}\\
磁化等离子体碰撞输运的权威参考书。详细推导了新经典输运（香蕉轨道、 Pfirsch--Schlüter 流、靴带电流），是理解托卡马克输运瓶颈的必读之作。

\item \textbf{Thomas H. Stix, \textit{Waves in Plasmas}, Springer.}\\
等离子体波理论的经典专著。对冷等离子体波、热波、动力学波（包括 Bernstein 波、漂移波）的色散关系、极化、阻尼给出了统一且系统的处理。阅读门槛较高，需熟悉复变函数与 Landau 围道。

\item \textbf{Paul M. Bellan, \textit{Fundamentals of Plasma Physics}, Cambridge.}\\
从电磁学基本方程出发，对 MHD、动力学、波、不稳定性有全面而现代的处理。特色是大量物理图像和数值估计，适合作为案头参考。
\end{enumerate}


\section{科研论文导读}

以下按方向推荐奠基性论文，标注阅读门槛。

\begin{table}[htbp]
\centering
\caption{经典科研论文推荐}
\label{tab:papers}
\small
\begin{tabular}{@{}p{2.5cm}p{5.5cm}p{5.5cm}@{}}
\toprule
\textbf{方向} & \textbf{论文} & \textbf{阅读门槛与价值} \\
\midrule
\multirow{2}{2.5cm}{磁约束聚变} & L. Spitzer, \textit{Physics of Fully Ionized Gases}, 1956 (Interscience). & 经典输运理论的奠基。门槛：中等，需熟悉 Fokker--Planck 方程。 \\
 & H. P. Furth \textit{et al.}, ``Tokamak equilibria with finite $\beta$'', Phys. Fluids 1970. & 高比压托卡马克平衡的里程碑。门槛：高，需 Grad--Shafranov 方程基础。 \\
\midrule
\multirow{2}{2.5cm}{空间等离子体} & J. W. Dungey, ``Interplanetary magnetic field and the auroral zones'', 1961. & 磁重联概念的奠基。门槛：低，物理图像极其清晰。 \\
 & A. Hasegawa, \textit{Plasma Instabilities and Nonlinear Effects}, 1975 (Springer). & 空间等离子体中漂移波、阿尔芬波不稳定性的经典教材。门槛：中等。 \\
\midrule
\multirow{2}{2.5cm}{低温/工业等离子体} & M. A. Lieberman \& A. J. Lichtenberg, \textit{Principles of Plasma Discharges and Materials Processing}, 2nd ed., Wiley. & 低温等离子体物理与工艺的标准教材。门槛：中等，需气体放电基础。 \\
 & L. Tonks \& I. Langmuir, ``Oscillations in ionized gases'', Phys. Rev. 1929. & 等离子体振荡的首次实验观测与理论。门槛：低，历史意义重大。 \\
\bottomrule
\end{tabular}
\end{table}


%======================================================================
% 附录F：全书高频易错点汇总
%======================================================================
\chapter*{附录F：全书高频易错点汇总}
\addcontentsline{toc}{chapter}{附录F：全书高频易错点汇总}

本章汇总全书（第1--12章）中读者最容易产生的误解。每条先给出 \textbf{常见错误说法}，再给出 \textbf{正确理解}。建议在学习对应章节时对照自查。


\section{第1--2章：等离子体基本概念}

\begin{warning}
\begin{enumerate}
\item \textbf{错误：}「准中性 = 完全无电场」\\
\textbf{正确：} 准中性意味着在远大于德拜长度的尺度上宏观电荷密度为零；但在德拜球内部，局部可以存在显著的净电荷和电场。

\item \textbf{错误：}「等离子体必须带电」\\
\textbf{正确：} 等离子体是 \textbf{准中性} 的——正负电荷几乎相等，整体不显示净电荷。它的特殊行为来自集体相互作用，而非净电荷。

\item \textbf{错误：}「$N_D\gg 1$ 意味着每个粒子都不受库仑力影响」\\
\textbf{正确：} $N_D\gg 1$ 意味着德拜屏蔽效应很强，远距离库仑力被屏蔽；但近距离碰撞（$b<b_\pi$）仍然频繁发生，碰撞输运不可忽略。
\end{enumerate}
\end{warning}


\section{第3--4章：单粒子运动与碰撞}

\begin{warning}
\begin{enumerate}
\setcounter{enumi}{3}
\item \textbf{错误：}「$\vect{E}\times\vect{B}$ 漂移与电荷 $q$ 有关」\\
\textbf{正确：} $\vect{v}_E = \vect{E}\times\vect{B}/B^2$ 与 $q$ 无关，电子和离子同向漂移。这是漂移运动中最易被混淆的结论。

\item \textbf{错误：}「梯度漂移 $\vect{v}_\nabla B$ 和曲率漂移 $\vect{v}_R$ 可以分别抵消」\\
\textbf{正确：} 在真空场 $\curl\vect{B}=0$ 中，两者大小相等、方向相同（对正离子），因此 \textbf{叠加} 而非抵消。总漂移为 $2\times$ 单一项。

\item \textbf{错误：}「极化漂移是稳态漂移」\\
\textbf{正确：} 极化漂移 $\vect{v}_p = (m/qB^2)\,\mathrm{d}\vect{E}_\perp/\mathrm{d}t$ 明确依赖于电场随时间的 \textbf{变化率}，仅存在于非稳态（启动、关断、波场）中。

\item \textbf{错误：}「碰撞频率与温度成正比」\\
\textbf{正确：} 对于高温库仑碰撞，$\nu\propto T^{-3/2}$——温度越高，碰撞越慢（因为粒子飞得越快，相互作用时间越短）。

\item \textbf{错误：}「库仑对数 $\ln\Lambda$ 可以随便取 10--20」\\
\textbf{正确：} $\ln\Lambda$ 虽变化缓慢，但在高密度、低温等离子体（如尘埃等离子体、惯性约束）中可能降至 $<5$，此时必须精确计算，不能随意取常数。
\end{enumerate}
\end{warning}


\section{第5--6章：流体方程与输运}

\begin{warning}
\begin{enumerate}
\setcounter{enumi}{7}
\item \textbf{错误：}「双极扩散中电子和离子扩散系数相同」\\
\textbf{正确：} 电子和离子扩散系数 $D_e, D_i$ 差异巨大；但电荷分离产生的内建电场迫使它们以共同的双极扩散系数 $D_a$ 运动。$D_a\approx D_i(1+T_e/T_i)$，被电子温度放大，但由离子扩散主导。

\item \textbf{错误：}「广义欧姆定律中霍尔项总是很小」\\
\textbf{正确：} 在强磁场、小尺度（如磁重联区、霍尔推进器）中，霍尔项 $-\vect{j}\times\vect{B}/(ne)$ 与电阻项相当甚至更大，不可忽略。

\item \textbf{错误：}「理想 MHD 意味着完全无电阻」\\
\textbf{正确：} 理想 MHD 假设 $R_m\gg 1$（磁冻结），但电阻 $\eta$ 仍然存在，只是对流项主导。在边界层、电流片中，$R_m\sim O(1)$，必须回到非理想方程。
\end{enumerate}
\end{warning}


\section{第7--8章：波与动力学}

\begin{warning}
\begin{enumerate}
\setcounter{enumi}{10}
\item \textbf{错误：}「Landau 阻尼是碰撞导致的能量耗散」\\
\textbf{正确：} Landau 阻尼是纯粹的 \textbf{无碰撞} 集体效应——波与共振粒子的能量交换。$dF_0/dv<0$ 时波能量被粒子吸收，$dF_0/dv>0$（如束流）时反而出现增长。

\item \textbf{错误：}「哨声波是声波的一种」\\
\textbf{正确：} 哨声波是 \textbf{电磁波}，右旋圆极化（R 波）的低频分支。它之所以叫 ``whistler''，是因为其群速度 $\diff\omega/\diff k$ 依赖于频率，导致不同频率分量到达时间不同，在音频上形成下降的哨音。

\item \textbf{错误：}「傅里叶变换的符号约定无关紧要」\\
\textbf{正确：} 本书统一使用 \textbf{物理学约定}：平面波形式为 $\exp[-\ii(\vect{k}\cdot\vect{r}-\omega t)]$。这与工程学中常用 $\exp[+\ii(\omega t - \vect{k}\cdot\vect{r})]$ 的约定不同，导致色散关系中虚部符号相反。请务必与本书保持一致，否则 Landau 阻尼会变号。

\item \textbf{错误：}「截止频率和传播频率是一回事」\\
\textbf{正确：} ``截止''（cutoff）意味着 $k\to 0$，波无法传播；``共振''（resonance）意味着 $k\to\infty$，波被强烈吸收。二者物理完全不同。
\end{enumerate}
\end{warning}


\section{第9--11章：MHD、平衡与不稳定性}

\begin{warning}
\begin{enumerate}
\setcounter{enumi}{14}
\item \textbf{错误：}「MHD 可以描述所有等离子体现象」\\
\textbf{正确：} MHD 忽略了霍尔效应、双极扩散、电子惯性、有限拉莫尔半径效应。在动力学尺度（如磁重联的扩散区、边缘等离子体）中，必须回到双流体或动力学理论。

\item \textbf{错误：}「磁冻结意味着磁力线永远不动」\\
\textbf{正确：} 磁冻结意味着等离子体流体元与磁通量一起运动；在理想 MHD 中，磁力线可随流体整体运动，但不会在流体内部 ``切断'' 或 ``重联''。

\item \textbf{错误：}「安全因子 $q>1$ 就足够稳定」\\
\textbf{正确：} $q>1$（Kruskal--Shafranov 极限）仅排除最低阶扭折（kink）模式；更高阶扭曲、气球、撕裂模等不稳定性还需要 $q>2$ 或更严格的剖面约束。$q$ 的径向分布（$q$ 剖面）比单一数值更重要。

\item \textbf{错误：}「比压 $\beta$ 越高越好」\\
\textbf{正确：} 高 $\beta$ 意味着等离子体压力占比大，能量利用率高，但 MHD 气球不稳定性在 $\beta$ 超过某个阈值（与 $q$、剖面形状有关）时被激发。聚变堆设计必须在 ``高 $\beta$ 追求'' 与 ``稳定性裕度'' 之间取折中。

\item \textbf{错误：}「磁重联只发生在 X 点」\\
\textbf{正确：} X 点（磁零点附近）是磁重联的核心区域，但重联是一个宏观过程，涉及入流区、出流区、扩散区的全局耦合。二维 X 点模型只是简化图像。
\end{enumerate}
\end{warning}


\section{第12章：应用与前沿}

\begin{warning}
\begin{enumerate}
\setcounter{enumi}{19}
\item \textbf{错误：}「托卡马克中所有等离子体电流都是外部驱动的」\\
\textbf{正确：} 托卡马克中存在显著的 \textbf{自举电流}（bootstrap current）——由碰撞导致的新经典输运驱动的内部电流，可占总电流的 $50\%$ 以上，对维持稳态运行至关重要。

\item \textbf{错误：}「空间等离子体总是无碰撞的」\\
\textbf{正确：} 太阳大气、电离层、行星磁层中，碰撞效应（如库仑碰撞、电荷交换）在特定区域（如磁层等离子体片、日冕底部）仍然重要，不能一概而论。
\end{enumerate}
\end{warning}


%======================================================================
% 附录G：Python仿真入门
%======================================================================
\chapter*{附录G：Python仿真入门}
\addcontentsline{toc}{chapter}{附录G：Python仿真入门}

本章为读者提供使用 Python 进行等离子体物理数值模拟的入门指导。代码使用 \texttt{NumPy}、\texttt{SciPy}、\texttt{Matplotlib} 等标准库，可直接运行。


\section{G.1 单粒子轨道追踪}

\textbf{物理目标：} 数值求解洛伦兹力方程
\begin{equation}
m\frac{\diff\vect{v}}{\diff t} = q(\vect{E}+\vect{v}\times\vect{B}), \qquad \frac{\diff\vect{r}}{\diff t} = \vect{v}
\end{equation}
演示均匀磁场中的回旋运动、$\vect{E}\times\vect{B}$ 漂移和磁镜反射。

\begin{tip}
运行前请确保已安装：\texttt{pip install numpy scipy matplotlib}
\end{tip}

\begin{verbatim}
import numpy as np
import matplotlib.pyplot as plt
from scipy.integrate import odeint

# -------------------- 物理参数 --------------------
q = 1.602e-19       # 电荷 (C)
m = 1.673e-27       # 质量 (kg, 质子)
B0 = np.array([0, 0, 1.0])  # 均匀磁场 (T)，沿 z
E0 = np.array([0, 1e3, 0])  # 电场 (V/m)，沿 y

# -------------------- 洛伦兹力方程 --------------------
def lorentz(y, t, q, m, E, B):
    """
    y = [x, y, z, vx, vy, vz]
    返回 dy/dt
    """
    r = y[0:3]
    v = y[3:6]
    # F = q(E + v x B)
    F = q * (E + np.cross(v, B))
    a = F / m
    return np.concatenate([v, a])

# -------------------- 初始条件 --------------------
r0 = np.array([0, 0, 0])          # 初始位置 (m)
v0 = np.array([1e5, 0, 1e4])      # 初始速度 (m/s)
y0 = np.concatenate([r0, v0])

# 时间网格
t = np.linspace(0, 1e-6, 2000)    # 1微秒，2000步

# 数值积分
sol = odeint(lorentz, y0, t, args=(q, m, E0, B0))

# 提取结果
x, y, z = sol[:, 0], sol[:, 1], sol[:, 2]
vx, vy, vz = sol[:, 3], sol[:, 4], sol[:, 5]

# -------------------- 可视化 --------------------
fig = plt.figure(figsize=(12, 5))

# 3D 轨迹
ax1 = fig.add_subplot(131, projection='3d')
ax1.plot(x, y, z, lw=0.8)
ax1.set_title("3D Orbit")
ax1.set_xlabel('x [m]'); ax1.set_ylabel('y [m]'); ax1.set_zlabel('z [m]')

# x-y 平面投影（回旋 + E x B 漂移）
ax2 = fig.add_subplot(132)
ax2.plot(x, y, lw=0.8)
ax2.set_aspect('equal')
ax2.set_title("x-y Projection")
ax2.set_xlabel('x [m]'); ax2.set_ylabel('y [m]')

# v_z 随时间变化（检查是否守恒）
ax3 = fig.add_subplot(133)
ax3.plot(t * 1e6, vz, lw=0.8)
ax3.set_title("v_z vs time")
ax3.set_xlabel('t [us]'); ax3.set_ylabel('v_z [m/s]')

plt.tight_layout()
plt.savefig('single_particle_orbit.png', dpi=150)
plt.show()
\end{verbatim}

\begin{insight}{}{}
在均匀磁场中，洛伦兹力方程可解析求解：粒子在垂直于 $B$ 的平面内做圆周运动（回旋），圆心以速度 $\vect{E}\times\vect{B}/B^2$ 漂移。若 $B$ 不均匀（如磁镜场），则 $v_\parallel$ 与 $v_\perp$ 之间发生交换，导致磁镜反射。读者可将上例中的 $B0$ 改为 $B(z)=B_0(1+(z/L)^2)$ 的形式，观察 $v_\parallel$ 在磁场增强区趋于零的现象。
\end{insight}


\section{G.2 德拜屏蔽电势的数值求解}

\textbf{物理目标：} 求解一维球坐标泊松方程
\begin{equation}
\frac{1}{r^2}\frac{\diff}{\diff r}\left(r^2\frac{\diff\phi}{\diff r}\right) = \frac{e n_0}{\eps_0}\left[\exp\!\left(\frac{e\phi}{k_\mathrm{B}T_e}\right)-1\right]
\end{equation}
与线性化德拜屏蔽解 $\phi(r)=\frac{e}{4\pi\eps_0 r}\exp(-r/\lambda_D)$ 对比。

\begin{verbatim}
import numpy as np
import matplotlib.pyplot as plt
from scipy.integrate import solve_bvp

# -------------------- 物理参数 --------------------
n0 = 1e18           # 背景密度 m^-3
Te = 10.0           # 电子温度 eV
q_e = 1.602e-19
eps0 = 8.854e-12
kB = 1.381e-23

# 德拜长度
lam_D = np.sqrt(eps0 * kB * (Te * 11604) / (n0 * q_e**2))

# -------------------- 定义ODE --------------------
def poisson_debye(r, y):
    """
    y = [phi, dphi/dr]
    方程: d^2phi/dr^2 + (2/r)dphi/dr = (q_e*n0/eps0)*(exp(q_e*phi/(kB*Te)) - 1)
    """
    phi, dphi = y
    # 为避免 r=0 奇点，从 rmin > 0 开始积分
    rhs = (q_e * n0 / eps0) * (np.exp(q_e * phi / (kB * Te * 11604)) - 1)
    if np.any(r == 0):
        rhs = 0  # 在原点处用极限值
    ddphi = rhs - (2.0 / r) * dphi
    return np.vstack([dphi, ddphi])

# 边界条件
# phi(inf) = 0, phi'(inf) = 0
# phi(rmin) = e/(4*pi*eps0*rmin)  (近似点电荷)
def bc(ya, yb):
    r_min = 1e-6
    phi_b = q_e / (4 * np.pi * eps0 * r_min)
    return np.array([ya[0] - phi_b, yb[0], yb[1]])

# 网格
r = np.linspace(1e-6, 5*lam_D, 500)
y_guess = np.zeros((2, r.size))

# 求解
sol = solve_bvp(poisson_debye, bc, r, y_guess, max_nodes=1e5, tol=1e-6)

# 解析解
phi_analytic = (q_e / (4 * np.pi * eps0 * r)) * np.exp(-r / lam_D)

# -------------------- 绘图 --------------------
plt.figure(figsize=(8, 5))
plt.semilogy(r / lam_D, np.abs(sol.y[0]), label='Nonlinear Poisson', lw=2)
plt.semilogy(r / lam_D, np.abs(phi_analytic), '--', label='Debye: $e^{-r/\\lambda_D}$', lw=2)
plt.xlabel('$r / \\lambda_D$')
plt.ylabel('$|\\phi|$ [V]')
plt.title('Debye Shielding Potential')
plt.legend()
plt.grid(True, which='both', ls='--', alpha=0.5)
plt.tight_layout()
plt.savefig('debye_shielding.png', dpi=150)
plt.show()
\end{verbatim}


\section{G.3 朗缪尔波色散曲线绘制}

\textbf{物理目标：} 绘制冷等离子体朗缪尔波 $\omega_{pe}$ 加上热修正（Bohm--Gross 色散关系）：
\begin{equation}
\omega^2 = \omega_{pe}^2 + \frac{3k_\mathrm{B}T_e}{m_e}k^2 = \omega_{pe}^2 + 3k^2v_{te}^2
\end{equation}

\begin{verbatim}
import numpy as np
import matplotlib.pyplot as plt

# -------------------- 物理参数 --------------------
n = 1e18            # m^-3
Te = 10.0           # eV
q_e = 1.602e-19
m_e = 9.109e-31
kB = 1.381e-23

# 等离子体频率
omega_pe = np.sqrt(n * q_e**2 / (8.854e-12 * m_e))
# 热速度
v_te = np.sqrt(kB * (Te * 11604) / m_e)

# 波数网格 (限制在 k*lambda_D < 1 内，Bohm-Gross近似有效)
lam_D = v_te / omega_pe
k = np.linspace(0.01, 1.0 / lam_D, 500)

# Bohm-Gross 色散关系
omega = np.sqrt(omega_pe**2 + 3 * k**2 * v_te**2)

# -------------------- 绘图 --------------------
plt.figure(figsize=(8, 5))
plt.plot(k * lam_D, omega / omega_pe, 'b-', lw=2, label='Bohm-Gross')
plt.axhline(y=1.0, color='r', linestyle='--', label='Cold plasma: $\\omega=\\omega_{pe}$')
plt.xlabel('$k \\lambda_D$')
plt.ylabel('$\\omega / \\omega_{pe}$')
plt.title('Langmuir Wave Dispersion Relation')
plt.legend()
plt.grid(True, ls='--', alpha=0.5)
plt.xlim(0, 1.0)
plt.ylim(0.9, 2.5)
plt.tight_layout()
plt.savefig('langmuir_dispersion.png', dpi=150)
plt.show()
\end{verbatim}


\section{G.4 仿真软件清单与学习路线}

本节列出等离子体物理领域主流的数值模拟代码，按方法分类，并给出学习路线建议。

\begin{table}[htbp]
\centering
\caption{等离子体数值模拟软件清单}
\label{tab:software}
\begin{tabular}{@{}p{2.8cm}p{3.5cm}p{6.5cm}@{}}
\toprule
\textbf{类型} & \textbf{软件名称} & \textbf{特点与适用场景} \\
\midrule
\multirow{3}{2.8cm}{PIC 粒子模拟} & EPOCH & 英国华威大学开发，开源，支持 1D/2D/3D，电磁和静电 PIC，文档完善，适合激光等离子体、加速器等。 \\
 & SMILEI & 法国开源，高可扩展性，支持 GPU 加速，负载均衡优秀，适合大规模并行。 \\
 & OSIRIS & UCLA 开发，相对论 PIC 标杆，在激光尾场加速、辐射物理中广泛使用。 \\
\midrule
\multirow{2}{2.8cm}{流体/ MHD} & BOUT++ & 英国开发，基于有限差分，专注边界等离子体、湍流、MHD 不稳定性，托卡马克偏滤器模拟首选。 \\
 & Nektar++ & 谱元法高阶精度，适合高雷诺数流动、阿尔芬波、流体力学问题。 \\
\midrule
\multirow{2}{2.8cm}{动理学/ 回旋动理学} & GYRO & 美国通用原子能开发，托卡马克湍流输运（ITG、TEM）的标准工具。 \\
 & GENE & 德国开发，基于谱方法，支持全局和局部模拟，广泛用于聚变动理学研究。 \\
\midrule
\multirow{2}{2.8cm}{磁重联/ 多流体} & P3D & 美国开发，3D PIC 和 MHD 混合，空间等离子体磁重联研究常用。 \\
 & Pluto & 意大利开发，高分辨率 Godunov 型 MHD/RMHD 代码，天体物理喷流、吸积盘适用。 \\
\bottomrule
\end{tabular}
\end{table}

\begin{tip}
\textbf{推荐学习路线：}
\begin{enumerate}
\item \textbf{阶段一（Python 练习）：} 先完成本书附录 G.1--G.3 的代码，理解单粒子轨道、泊松方程、色散关系的数值实现。
\item \textbf{阶段二（1D PIC）：} 使用 EPOCH 的 1D 教程，运行朗缪尔波、双流不稳定性（two-stream instability）的基准算例。
\item \textbf{阶段三（2D/3D）：} 在集群上运行自己的 2D 问题（如磁重联、Landau 阻尼），学习并行 I/O、后处理（VisIt / ParaView）。
\item \textbf{阶段四（深入研究）：} 根据方向选择：聚变输运 $\to$ GYRO/GENE；空间物理 $\to$ P3D/SMILEI；工业低温 $\to$ BOLSIG+ / COMSOL Plasma Module。
\end{enumerate}
\end{tip}

% 附录结束

